\documentclass[11pt]{ctexart}
\usepackage[a4paper,margin=1in]{geometry}
\usepackage{longtable,booktabs}
\title{Geant4-Agent 回归测试报告（中文）}
\author{自动化评测流程}
\date{2026-03-08\_124606}
\begin{document}
\maketitle
\section{统一 Pass/Fail 总表}
\begin{longtable}{p{0.06\linewidth}p{0.08\linewidth}p{0.07\linewidth}p{0.08\linewidth}p{0.07\linewidth}p{0.07\linewidth}p{0.07\linewidth}p{0.07\linewidth}p{0.09\linewidth}p{0.06\linewidth}p{0.21\linewidth}}
\toprule
ID & 领域 & 风格 & 参数完整性 & 期望初始 & 实际初始 & 期望最终 & 实际最终 & 缺参(初->终) & 结果 & 失败原因 \\
\midrule
MTM01 & medical & fuzzy & missing & 失败 & 失败 & 通过 & 通过 & 12->0 & 通过 & - \\
MTM02 & aerospace & precise & missing & 失败 & 失败 & 通过 & 通过 & 2->0 & 通过 & - \\
MTM03 & industrial\_ndt & precise & missing & 失败 & 失败 & 通过 & 通过 & 12->0 & 通过 & - \\
MTM04 & physics & precise & missing & 失败 & 失败 & 通过 & 通过 & 8->0 & 通过 & - \\
MTM05 & nuclear\_engineering & fuzzy & missing & 失败 & 失败 & 通过 & 通过 & 5->0 & 通过 & - \\
MTM06 & security\_screening & fuzzy & missing & 失败 & 失败 & 通过 & 通过 & 11->0 & 通过 & - \\
MTM07 & semiconductor & precise & missing & 失败 & 失败 & 通过 & 通过 & 3->0 & 通过 & - \\
MTM08 & education & fuzzy & missing & 失败 & 失败 & 通过 & 通过 & 11->0 & 通过 & - \\
MTM09 & environmental\_radiation & fuzzy & missing & 失败 & 失败 & 通过 & 通过 & 3->0 & 通过 & - \\
MTM10 & high\_energy\_physics & precise & missing & 失败 & 失败 & 通过 & 通过 & 10->0 & 通过 & - \\
MTM11 & materials\_update & precise & full & 通过 & 通过 & 通过 & 通过 & 0->0 & 通过 & - \\
MTM12 & overwrite\_reject & precise & full & 通过 & 通过 & 通过 & 通过 & 0->0 & 通过 & - \\
MTM13 & medical\_detector & fuzzy & missing & 失败 & 失败 & 通过 & 通过 & 13->0 & 通过 & - \\
MTM14 & industrial\_panel & precise & missing & 失败 & 失败 & 通过 & 通过 & 15->0 & 通过 & - \\
MTM15 & shielding\_stack & precise & missing & 失败 & 失败 & 通过 & 通过 & 14->0 & 通过 & - \\
MTM16 & nested\_target & precise & missing & 失败 & 失败 & 通过 & 通过 & 13->0 & 通过 & - \\
MTM17 & coaxial\_shield & fuzzy & missing & 失败 & 失败 & 通过 & 通过 & 10->0 & 通过 & - \\
MTM18 & boolean\_cavity & precise & missing & 失败 & 失败 & 通过 & 通过 & 17->0 & 通过 & - \\
MTM19 & ring\_complete & precise & full & 通过 & 通过 & 通过 & 通过 & 0->0 & 通过 & - \\
MTM20 & grid\_complete & precise & full & 通过 & 通过 & 通过 & 通过 & 0->0 & 通过 & - \\
MTM21 & layered\_detector & fuzzy & missing & 失败 & 失败 & 通过 & 通过 & 14->0 & 通过 & - \\
MTM22 & embedded\_sensor & precise & full & 通过 & 通过 & 通过 & 通过 & 0->0 & 通过 & - \\
MTM23 & layered\_shield\_complete & precise & full & 通过 & 通过 & 通过 & 通过 & 0->0 & 通过 & - \\
MTM24 & boolean\_complete & precise & full & 通过 & 通过 & 通过 & 通过 & 0->0 & 通过 & - \\
\bottomrule
\end{longtable}
\section{基准对比}
\begin{itemize}
\item pass\_rate=1.000
\item initial\_alignment\_rate=1.000
\item final\_alignment\_rate=1.000
\item avg\_turns\_per\_case=2.458
\item missing\_reduction\_rate=1.000
\item confirm\_apply\_success\_rate=1.000
\item rollback\_after\_confirm\_rate=0.000
\item internal\_leak\_turn\_rate=0.068
\item repeated\_assistant\_case\_rate=0.000
\item stale\_confirm\_case\_rate=0.000
\item delta\_pass\_rate=+0.050
\item regressions=无
\item improvements=无
\end{itemize}
\appendix
\section{附录：完整对话与内部处理链（自然语言）}
\subsection{案例 MTM01（medical）}
\textbf{对话：}
\begin{itemize}
\item 用户：我先要一个用于医学成像的初步辐照场景，细节还没定，你先告诉我还缺什么。
\item 系统：目前我们已经确定了使用盒体作为几何结构，接下来还需要补充一些具体信息：

- **几何**：需要盒体在x、y、z三个方向上的具体尺寸。
- **材料**：需要指定体积与材料的绑定关系。
- **源**：需要选择源的类型（例如粒子束、放射源等）。

请随时告诉我这些细节，我们可以一步步完善这个辐照场景。
\item 用户：采用 300 mm x 300 mm x 300 mm 的水箱，材料 ?；源用束流 gamma，能量 6 MeV，位置 (0,0,-400) mm，方向 +z；物理列表 QBBC；输出 root。
\item 系统：我检测到这会覆盖已经确认的值。待修改内容：几何结构类型：unknown -> single\_box；盒体 x 方向尺寸：None -> 300.0。如果要应用这些修改，请直接回复“确认”；如果要保留原值，请直接说“保持原值”。
\item 用户：确认
\item 系统：配置已完成。当前已确认的几何参数包括：盒体在 x、y、z 方向上的尺寸以及几何结构类型。
\end{itemize}
\textbf{内部链路（自然语言）：}
\begin{itemize}
\item 第 1 轮：动作=summarize\_progress，阶段=geometry，闭环状态=False。
\item 缺失字段：geometry.params.module\_x、geometry.params.module\_y、geometry.params.module\_z、materials.volume\_material\_map、source.type、source.particle、source.energy、source.position、source.direction、physics.physics\_list、output.format、output.path。
\item LLM 使用=True；回退=无；后端=llm\_semantic\_frame+runtime\_semantic。
\item 第 2 轮：动作=confirm\_overwrite，阶段=geometry，闭环状态=False。
\item 缺失字段：geometry.params.module\_x、geometry.params.module\_y、geometry.params.module\_z。
\item LLM 使用=True；回退=无；后端=llm\_slot\_frame+runtime\_semantic。
\item 第 3 轮：动作=finalize，阶段=materials，闭环状态=True。
\item 缺失字段：无。
\item LLM 使用=True；回退=无；后端=llm\_semantic\_frame+runtime\_semantic+confirmed\_pending\_overwrite。
\end{itemize}
\subsection{案例 MTM02（aerospace）}
\textbf{对话：}
\begin{itemize}
\item 用户：I need a first-pass proton irradiation setup for a silicon device; keep it simple for now.
\item 系统：Geometry: updated geometry structure, box size along x, box size along y. Materials: updated selected materials, volume-to-material binding. Source: updated source type, particle type, source energy. Physics: updated physics list. Output: still needs output format, output path.
\item 用户：Use a silicon box 50 mm x 50 mm x 5 mm, beam source proton 100 MeV from (0,0,-100) mm along +z, physics FTFP\_BERT, output json.
\item 系统：You're about to overwrite a previously confirmed value.

Pending change:
- Box size along x: 1000.0 → 50.0
- Box size along y: 1000.0 → 50.0

Reply **'confirm'** to apply these changes, or **'keep existing'** to keep the current values.
\item 用户：confirm
\item 系统：Configuration complete. Confirmed items: source type, box size along z, box size along y, and box size along x.
\end{itemize}
\textbf{内部链路（自然语言）：}
\begin{itemize}
\item 第 1 轮：动作=summarize\_progress，阶段=materials，闭环状态=False。
\item 缺失字段：output.format、output.path。
\item LLM 使用=True；回退=无；后端=llm\_slot\_frame+runtime\_semantic。
\item 第 2 轮：动作=confirm\_overwrite，阶段=source，闭环状态=False。
\item 缺失字段：无。
\item LLM 使用=True；回退=无；后端=llm\_slot\_frame+runtime\_semantic。
\item 第 3 轮：动作=finalize，阶段=physics，闭环状态=True。
\item 缺失字段：无。
\item LLM 使用=True；回退=无；后端=llm\_semantic\_frame+runtime\_semantic+confirmed\_pending\_overwrite。
\end{itemize}
\subsection{案例 MTM03（industrial\_ndt）}
\textbf{对话：}
\begin{itemize}
\item 用户：我要做一个 gamma 穿透铁靶的小测试，但我还没把完整参数说完。
\item 系统：目前我们已确定使用铁作为材料，并选择 gamma 粒子作为源类型。

接下来还需要你补充几何部分的信息，包括结构类型、盒体在 x、y、z 方向的尺寸，以及将体积与材料绑定的设置。
\item 用户：请设为半径 40 mm、半长 80 mm 的铁圆柱，材料 ?；点源 gamma，能量 2 MeV，位置 (0,0,-200) mm，方向 +z；物理列表 Shielding；输出 csv。
\item 系统：配置已完成！当前已确认的设置包括：体积与材料绑定、输出路径、输出格式和物理列表。
\end{itemize}
\textbf{内部链路（自然语言）：}
\begin{itemize}
\item 第 1 轮：动作=summarize\_progress，阶段=geometry，闭环状态=False。
\item 缺失字段：geometry.structure、geometry.params.module\_x、geometry.params.module\_y、geometry.params.module\_z、materials.volume\_material\_map、source.type、source.energy、source.position、source.direction、physics.physics\_list、output.format、output.path。
\item LLM 使用=True；回退=无；后端=llm\_slot\_frame+runtime\_semantic。
\item 第 2 轮：动作=finalize，阶段=materials，闭环状态=True。
\item 缺失字段：无。
\item LLM 使用=True；回退=无；后端=llm\_slot\_frame+runtime\_semantic。
\end{itemize}
\subsection{案例 MTM04（physics）}
\textbf{对话：}
\begin{itemize}
\item 用户：Gamma point source 1 MeV at (0,0,-100) toward +z; please build the rest.
\item 系统：Geometry: still needs geometry structure, box size along x. Materials: still needs volume-to-material binding. Source: updated source type, particle type, source energy. Physics: still needs physics list. Output: still needs output format, output path.
\item 用户：Set copper box 1 m x 1 m x 1 m, physics FTFP\_BERT, output hdf5.
\item 系统：Your configuration is complete. The following items have been confirmed: volume-to-material binding, output path, output format, and physics list.
\end{itemize}
\textbf{内部链路（自然语言）：}
\begin{itemize}
\item 第 1 轮：动作=summarize\_progress，阶段=geometry，闭环状态=False。
\item 缺失字段：geometry.structure、geometry.params.module\_x、geometry.params.module\_y、geometry.params.module\_z、materials.volume\_material\_map、physics.physics\_list、output.format、output.path。
\item LLM 使用=True；回退=无；后端=llm\_slot\_frame+runtime\_semantic。
\item 第 2 轮：动作=finalize，阶段=materials，闭环状态=True。
\item 缺失字段：无。
\item LLM 使用=True；回退=无；后端=llm\_slot\_frame+runtime\_semantic。
\end{itemize}
\subsection{案例 MTM05（nuclear\_engineering）}
\textbf{对话：}
\begin{itemize}
\item 用户：我想做一个中子屏蔽的概念验证，现在先给一个缺参开头。
\item 系统：目前我们已经确定了几个关键部分：几何结构、材料、粒子源类型、物理列表以及输出设置。

接下来还需要补充一些具体参数：盒体在x、y、z三个方向上的尺寸，以及源的发射能量和位置。

请随时告诉我这些信息，我们可以继续推进配置。
\item 用户：使用 200 mm x 200 mm x 200 mm 的铅块，材料 ?；点源 neutron，能量 14 MeV，位置 (0,0,-50) mm，方向 +z；物理列表 QGSP\_BERT\_HP；输出 xml。
\item 系统：我检测到这会覆盖已经确认的值。待修改内容：几何结构类型：unknown -> single\_box；已选材料：['G4\_Cu'] -> ['G4\_Pb']。如果要应用这些修改，请直接回复“确认”；如果要保留原值，请直接说“保持原值”。
\item 用户：确认
\item 系统：配置已完成。当前已确认的参数包括：体积与材料绑定、输出路径、盒体 z 方向尺寸、盒体 y 方向尺寸。
\end{itemize}
\textbf{内部链路（自然语言）：}
\begin{itemize}
\item 第 1 轮：动作=summarize\_progress，阶段=geometry，闭环状态=False。
\item 缺失字段：geometry.params.module\_x、geometry.params.module\_y、geometry.params.module\_z、source.energy、source.position。
\item LLM 使用=True；回退=无；后端=llm\_slot\_frame+runtime\_semantic。
\item 第 2 轮：动作=confirm\_overwrite，阶段=geometry，闭环状态=False。
\item 缺失字段：geometry.params.module\_x、geometry.params.module\_y、geometry.params.module\_z。
\item LLM 使用=True；回退=无；后端=llm\_slot\_frame+runtime\_semantic。
\item 第 3 轮：动作=finalize，阶段=materials，闭环状态=True。
\item 缺失字段：无。
\item LLM 使用=True；回退=无；后端=llm\_slot\_frame+runtime\_semantic+confirmed\_pending\_overwrite。
\end{itemize}
\subsection{案例 MTM06（security\_screening）}
\textbf{对话：}
\begin{itemize}
\item 用户：Need a quick cargo-screening toy setup with forward photons, but I have not fixed the target shape, material, source energy, or output format yet.
\item 系统：Geometry: updated geometry structure; still needs box size along x, box size along y. Materials: still needs volume-to-material binding. Source: updated particle type; still needs source type, source energy. Physics: still needs physics list. Output: still needs output format, output path.
\item 用户：Use steel cylinder target radius 60 mm half-length 120 mm, beam gamma 3 MeV from (0,0,-500) mm to +z, physics Shielding, output root.
\item 系统：This would overwrite an already confirmed value. Pending change: geometry structure: unknown -> single\_tubs; box size along x: None -> 1000.0. Reply 'confirm' to apply it, or say 'keep existing' to keep the current value.
\item 用户：confirm
\item 系统：Configuration complete. Confirmed items: volume-to-material binding, half length, radius, and box size along z.
\end{itemize}
\textbf{内部链路（自然语言）：}
\begin{itemize}
\item 第 1 轮：动作=summarize\_progress，阶段=geometry，闭环状态=False。
\item 缺失字段：geometry.params.module\_x、geometry.params.module\_y、geometry.params.module\_z、materials.volume\_material\_map、source.type、source.energy、source.position、source.direction、physics.physics\_list、output.format、output.path。
\item LLM 使用=True；回退=无；后端=llm\_slot\_frame+runtime\_semantic。
\item 第 2 轮：动作=confirm\_overwrite，阶段=geometry，闭环状态=False。
\item 缺失字段：geometry.params.module\_x、geometry.params.module\_y、geometry.params.module\_z。
\item LLM 使用=True；回退=无；后端=llm\_slot\_frame+runtime\_semantic。
\item 第 3 轮：动作=finalize，阶段=materials，闭环状态=True。
\item 缺失字段：无。
\item LLM 使用=True；回退=无；后端=llm\_semantic\_frame+runtime\_semantic+confirmed\_pending\_overwrite。
\end{itemize}
\subsection{案例 MTM07（semiconductor）}
\textbf{对话：}
\begin{itemize}
\item 用户：我想搭一个硅薄片的中子照射测试，但先只给一个简述。
\item 系统：目前我们已经完成了几何结构（盒体尺寸）、材料（硅）和粒子源（中子能量与位置）的设置。

接下来还需要你补充以下信息：
- 物理过程列表
- 输出文件格式
- 输出文件保存路径

请告诉我这些内容，我们就可以继续了。
\item 用户：使用 20 mm x 20 mm x 1 mm 的硅片，材料 ?；各向同性 neutron，能量 2 MeV，位置 (0,0,0) mm；物理列表 QGSP\_BERT\_HP；输出 json。
\item 系统：我注意到这次修改会覆盖已经确认的参数。具体变动如下：

- 盒体 x 方向尺寸：1000.0 → 20.0
- 盒体 y 方向尺寸：1000.0 → 20.0

请确认你的选择：
• 如果要应用这些修改，请直接回复“确认”。
• 如果要保留原来的数值，请直接回复“保持原值”。
\item 用户：确认
\item 系统：配置已完成。当前已确认的参数包括：源位置、源能量、源类型以及盒体 z 方向尺寸。
\end{itemize}
\textbf{内部链路（自然语言）：}
\begin{itemize}
\item 第 1 轮：动作=summarize\_progress，阶段=materials，闭环状态=False。
\item 缺失字段：physics.physics\_list、output.format、output.path。
\item LLM 使用=True；回退=无；后端=llm\_slot\_frame+runtime\_semantic。
\item 第 2 轮：动作=confirm\_overwrite，阶段=source，闭环状态=False。
\item 缺失字段：无。
\item LLM 使用=True；回退=无；后端=llm\_slot\_frame+runtime\_semantic。
\item 第 3 轮：动作=finalize，阶段=physics，闭环状态=True。
\item 缺失字段：无。
\item LLM 使用=True；回退=无；后端=llm\_semantic\_frame+runtime\_semantic+confirmed\_pending\_overwrite。
\end{itemize}
\subsection{案例 MTM08（education）}
\textbf{对话：}
\begin{itemize}
\item 用户：For classroom demo, I want gamma attenuation in copper; you can ask only essential missing parts.
\item 系统：Geometry: still needs geometry structure, box size along x. Materials: updated selected materials; still needs volume-to-material binding. Source: updated particle type; still needs source type, source energy. Physics: updated physics list. Output: still needs output format, output path.
\item 用户：Okay, copper box 100 mm x 100 mm x 100 mm, point gamma 1 MeV at (0,0,-200) mm, direction +z, physics FTFP\_BERT, output csv.
\item 系统：You've already confirmed the physics list as "Shielding." Changing it to "FTFP\_BERT" would overwrite that.

Please reply **'confirm'** to apply the change to FTFP\_BERT, or **'keep existing'** to keep Shielding.
\item 用户：confirm
\item 系统：Configuration complete. The following items are now confirmed: physics list, output path, output format, and source direction.
\end{itemize}
\textbf{内部链路（自然语言）：}
\begin{itemize}
\item 第 1 轮：动作=summarize\_progress，阶段=geometry，闭环状态=False。
\item 缺失字段：geometry.structure、geometry.params.module\_x、geometry.params.module\_y、geometry.params.module\_z、materials.volume\_material\_map、source.type、source.energy、source.position、source.direction、output.format、output.path。
\item LLM 使用=True；回退=无；后端=llm\_slot\_frame+runtime\_semantic。
\item 第 2 轮：动作=confirm\_overwrite，阶段=materials，闭环状态=False。
\item 缺失字段：无。
\item LLM 使用=True；回退=无；后端=llm\_slot\_frame+runtime\_semantic。
\item 第 3 轮：动作=finalize，阶段=source，闭环状态=True。
\item 缺失字段：无。
\item LLM 使用=True；回退=无；后端=llm\_semantic\_frame+runtime\_semantic+confirmed\_pending\_overwrite。
\end{itemize}
\subsection{案例 MTM09（environmental\_radiation）}
\textbf{对话：}
\begin{itemize}
\item 用户：我想先搭一个空气体积中的低能 gamma 传播场景，细节稍后补。
\item 系统：目前已经完成了几何、材料和源的基本设置。接下来还需要你补充物理列表、输出格式和输出路径这三项信息。
\item 用户：设为 500 mm x 500 mm x 500 mm 的空气盒，材料 ??；点源 gamma，能量 0.5 MeV，位置 (0,0,-300) mm，方向 +z；物理列表 QBBC；输出 root。
\item 系统：我注意到这次修改会覆盖你已经确认过的参数。具体变动如下：

- 盒体 x 方向尺寸：从 1000.0 改为 500.0
- 盒体 y 方向尺寸：从 1000.0 改为 500.0

请确认你的选择：
如果应用这些修改，请回复“确认”；
如果想保留原来的数值，请回复“保持原值”。
\item 用户：确认
\item 系统：配置已完成。当前已确认的参数包括：源位置、源能量、盒体 z 方向尺寸、盒体 y 方向尺寸。
\end{itemize}
\textbf{内部链路（自然语言）：}
\begin{itemize}
\item 第 1 轮：动作=summarize\_progress，阶段=materials，闭环状态=False。
\item 缺失字段：physics.physics\_list、output.format、output.path。
\item LLM 使用=True；回退=无；后端=llm\_slot\_frame+runtime\_semantic。
\item 第 2 轮：动作=confirm\_overwrite，阶段=source，闭环状态=False。
\item 缺失字段：无。
\item LLM 使用=True；回退=无；后端=llm\_slot\_frame+runtime\_semantic。
\item 第 3 轮：动作=finalize，阶段=physics，闭环状态=True。
\item 缺失字段：无。
\item LLM 使用=True；回退=无；后端=llm\_semantic\_frame+runtime\_semantic+confirmed\_pending\_overwrite。
\end{itemize}
\subsection{案例 MTM10（high\_energy\_physics）}
\textbf{对话：}
\begin{itemize}
\item 用户：Set geometry as a sphere target; I have not decided radius or source details yet.
\item 系统：Here's what we have so far:

**Geometry:** You've set the structure to a sphere. We still need the sphere's **radius**.

**Materials:** We need to assign a material to your volume.

**Source:** We need the source type, particle type, and particle energy.

**Physics \& Output:** We also need to select a physics list and set the output format and file path.
\item 用户：Sphere target radius 80 mm, lead material, electron point source 5 MeV at (0,0,-200) mm toward +z, physics FTFP\_BERT, output hdf5.
\item 系统：Your configuration is complete. The following items are confirmed: volume-to-material binding, output path, output format, and physics list.
\end{itemize}
\textbf{内部链路（自然语言）：}
\begin{itemize}
\item 第 1 轮：动作=summarize\_progress，阶段=geometry，闭环状态=False。
\item 缺失字段：geometry.params.child\_rmax、materials.volume\_material\_map、source.type、source.particle、source.energy、source.position、source.direction、physics.physics\_list、output.format、output.path。
\item LLM 使用=True；回退=无；后端=llm\_slot\_frame+runtime\_semantic。
\item 第 2 轮：动作=finalize，阶段=materials，闭环状态=True。
\item 缺失字段：无。
\item LLM 使用=True；回退=无；后端=llm\_slot\_frame+runtime\_semantic。
\end{itemize}
\subsection{案例 MTM11（materials\_update）}
\textbf{对话：}
\begin{itemize}
\item 用户：Create a complete setup: 100 mm x 100 mm x 100 mm copper box, point gamma 1 MeV at (0,0,-200) mm toward +z, physics QBBC, output json.
\item 系统：Current committed config is: physics list FTFP\_BERT, output format json, source type point, particle gamma. If you want changes, please specify which field to update.
\item 用户：Change the target material to aluminum and change the output format to root.
\item 系统：You're about to overwrite two confirmed settings:

- **Selected materials**: Copper → Aluminum
- **Output format**: JSON → ROOT

Please reply **"confirm"** to apply these changes, or **"keep existing"** to leave the current values as they are.
\item 用户：confirm
\item 系统：Configuration complete. The following items are confirmed: your selected materials, output format, output path, and volume-to-material bindings.
\end{itemize}
\textbf{内部链路（自然语言）：}
\begin{itemize}
\item 第 1 轮：动作=finalize，阶段=materials，闭环状态=True。
\item 缺失字段：无。
\item LLM 使用=True；回退=无；后端=llm\_slot\_frame+runtime\_semantic。
\item 第 2 轮：动作=confirm\_overwrite，阶段=source，闭环状态=False。
\item 缺失字段：无。
\item LLM 使用=True；回退=无；后端=llm\_slot\_frame+runtime\_semantic。
\item 第 3 轮：动作=finalize，阶段=physics，闭环状态=True。
\item 缺失字段：无。
\item LLM 使用=True；回退=无；后端=llm\_semantic\_frame+runtime\_semantic+confirmed\_pending\_overwrite。
\end{itemize}
\subsection{案例 MTM12（overwrite\_reject）}
\textbf{对话：}
\begin{itemize}
\item 用户：请创建一个完整配置：100 mm x 100 mm x 100 mm 的铜盒，材料 ?；点源 gamma，能量 1 MeV，位置 (0,0,-200) mm，方向 +z；物理列表 QBBC；输出 json。
\item 系统：配置已完成。当前已确认的设置包括：体积与材料绑定、输出路径和输出格式。
\item 用户：把材料改成 ?，并把输出格式改成 root。
\item 系统：检测到您想要修改输出格式，这将会覆盖已确认的当前值（json 改为 root）。请确认您的选择：直接回复“确认”以应用修改，或回复“保持原值”以保留原来的设置。
\item 用户：保持原值
\item 系统：好的，输出格式仍保持为 json，没有改动。
\end{itemize}
\textbf{内部链路（自然语言）：}
\begin{itemize}
\item 第 1 轮：动作=finalize，阶段=materials，闭环状态=True。
\item 缺失字段：无。
\item LLM 使用=True；回退=无；后端=llm\_slot\_frame+runtime\_semantic。
\item 第 2 轮：动作=confirm\_overwrite，阶段=source，闭环状态=False。
\item 缺失字段：geometry.ask.boolean.solid\_a\_size、geometry.ask.boolean.solid\_b\_size。
\item LLM 使用=True；回退=无；后端=llm\_slot\_frame+runtime\_semantic。
\item 第 3 轮：动作=reject\_overwrite，阶段=physics，闭环状态=True。
\item 缺失字段：无。
\item LLM 使用=True；回退=无；后端=llm\_semantic\_frame+runtime\_semantic+rejected\_pending\_overwrite。
\end{itemize}
\subsection{案例 MTM13（medical\_detector）}
\textbf{对话：}
\begin{itemize}
\item 用户：I want a ring of detector modules around a small source region, but I have not fixed the ring radius or module size yet.
\item 系统：Your geometry structure is set. Next, we need to specify the materials for your volumes, define your particle source, and choose the physics and output settings.
\item 用户：Use a ring of 12 modules, each 8 mm x 10 mm x 2 mm, radius 40 mm, clearance 1 mm, material water; point gamma source 0.662 MeV at (0,0,0) mm; physics QBBC; output json.
\item 系统：You are about to overwrite a confirmed geometry configuration.

**Current Geometry:**
*   Module: Box with dimensions 7.83 mm x 4.72 mm x 6.86 mm
*   Ring: 4 modules, radius 67.49 mm, clearance 0.19 mm

**New Geometry:**
*   Module: Box with dimensions 8.0 mm x 10.0 mm x 2.0 mm
*   Ring: 12 modules, radius 40.0 mm, clearance 1.0 mm

Please reply **'confirm'** to apply this change, or **'keep existing'** to keep the current values.
\item 用户：confirm
\item 系统：Configuration complete. The following geometry items are confirmed: ring clearance, ring radius, ring module count, and ring module size.
\end{itemize}
\textbf{内部链路（自然语言）：}
\begin{itemize}
\item 第 1 轮：动作=summarize\_progress，阶段=materials，闭环状态=False。
\item 缺失字段：materials.volume\_material\_map、source.type、source.particle、source.energy、source.position、source.direction、physics.physics\_list、output.format、output.path、geometry.ask.ring.module\_size、geometry.ask.ring.count、geometry.ask.ring.radius、geometry.ask.ring.clearance。
\item LLM 使用=True；回退=无；后端=llm\_semantic\_frame+runtime\_semantic。
\item 第 2 轮：动作=confirm\_overwrite，阶段=source，闭环状态=False。
\item 缺失字段：无。
\item LLM 使用=True；回退=无；后端=llm\_slot\_frame+runtime\_semantic。
\item 第 3 轮：动作=finalize，阶段=physics，闭环状态=True。
\item 缺失字段：无。
\item LLM 使用=True；回退=无；后端=llm\_semantic\_frame+runtime\_semantic+confirmed\_pending\_overwrite。
\end{itemize}
\subsection{案例 MTM14（industrial\_panel）}
\textbf{对话：}
\begin{itemize}
\item 用户：我想做一个二维阵列探测板，但 pitch 和材料还没定。
\item 系统：好的，目前配置进度如下：

**已完成：**
*   **几何**：已确定使用“二维阵列”结构。

**待补充：**
*   **几何**：阵列模块的尺寸、X方向的阵列数量。
*   **材料**：体积与材料的绑定关系。
*   **源**：源类型、粒子类型、能量和位置。
*   **物理**：物理过程列表。
*   **输出**：输出文件的格式和保存路径。

接下来，我们首先来确定阵列模块的具体尺寸吧。
\item 用户：设成 4 x 3 的盒体阵列，每个模块 10 mm x 10 mm x 2 mm，pitch\_x 12 mm，pitch\_y 14 mm，间隙 1 mm，材料 硅；点源 gamma 0.511 MeV，位置 (0,0,-80) mm，方向 +z；物理列表 QBBC；输出 root。
\item 系统：我注意到这将会覆盖已确认的几何参数。具体修改如下：

*   模块尺寸从 (7.83, 4.72, 6.86) mm 改为 (10.0, 10.0, 2.0) mm。
*   X方向间距从 11.47 mm 改为 12.0 mm。
*   Y方向间距从 8.35 mm 改为 14.0 mm。
*   阵列间隙从 1.82 mm 改为 1.0 mm。

请确认：**回复“确认”将应用以上修改；回复“保持原值”则保留现有配置。**
\item 用户：确认
\item 系统：配置已完成。当前已确认的参数包括：y方向阵列数量、x方向阵列数量、阵列间隙以及y方向的pitch。
\end{itemize}
\textbf{内部链路（自然语言）：}
\begin{itemize}
\item 第 1 轮：动作=summarize\_progress，阶段=materials，闭环状态=False。
\item 缺失字段：materials.volume\_material\_map、source.type、source.particle、source.energy、source.position、source.direction、physics.physics\_list、output.format、output.path、geometry.ask.grid.module\_size、geometry.ask.grid.count\_x、geometry.ask.grid.count\_y、geometry.ask.grid.pitch\_x、geometry.ask.grid.pitch\_y、geometry.ask.grid.clearance。
\item LLM 使用=True；回退=无；后端=llm\_semantic\_frame+runtime\_semantic。
\item 第 2 轮：动作=confirm\_overwrite，阶段=source，闭环状态=False。
\item 缺失字段：geometry.ask.grid.count\_x、geometry.ask.grid.count\_y。
\item LLM 使用=True；回退=无；后端=llm\_slot\_frame+runtime\_semantic。
\item 第 3 轮：动作=finalize，阶段=physics，闭环状态=True。
\item 缺失字段：无。
\item LLM 使用=True；回退=无；后端=llm\_semantic\_frame+runtime\_semantic+confirmed\_pending\_overwrite。
\end{itemize}
\subsection{案例 MTM15（shielding\_stack）}
\textbf{对话：}
\begin{itemize}
\item 用户：Build a stacked absorber sandwich along z; I will provide the layer sizes next.
\item 系统：You've set up a stacked absorber sandwich geometry along the z-axis. Here's what's complete and what's still needed:

**Completed:**
*   **Geometry:** The overall structure is set.

**Still Needed:**
*   **Geometry:** The footprint and thicknesses for each stack layer.
*   **Materials:** Which material to assign to each volume.
*   **Source:** The source type, particle type, and its energy.
*   **Physics:** The physics list to use.
*   **Output:** The format and file path for your results.
\item 用户：Use stacked layers with footprint 20 mm x 20 mm, thicknesses 2 mm, 4 mm, 6 mm, layer clearance 0.5 mm, outer box 30 mm x 30 mm x 20 mm, container clearance 1 mm, material aluminum; beam gamma 1.25 MeV from (0,0,-120) mm along +z; physics Shielding; output csv.
\item 系统：You are about to overwrite a confirmed geometry configuration.

**Current Geometry:**
*   **Stack:** Footprint \textasciitilde{}16.33 mm x 10.28 mm, Thicknesses [\textasciitilde{}6.68, \textasciitilde{}1.19, \textasciitilde{}5.59] mm, Clearance \textasciitilde{}0.73 mm.
*   **Container Box:** \textasciitilde{}18.07 mm x \textasciitilde{}45.52 mm x \textasciitilde{}16.66 mm, Clearance \textasciitilde{}0.87 mm.

**New Geometry:**
*   **Stack:** Footprint 20.0 mm x 20.0 mm, Thicknesses [2.0, 4.0, 6.0] mm, Clearance 0.5 mm.
*   **Container Box:** 30.0 mm x 30.0 mm x 20.0 mm, Clearance 1.0 mm.

Reply **'confirm'** to apply this change, or **'keep existing'** to keep the current values.
\item 用户：confirm
\item 系统：Your configuration is complete. The following geometry items are confirmed: stack container clearance, stack container size, stack layer clearance, and stack layer thicknesses.
\end{itemize}
\textbf{内部链路（自然语言）：}
\begin{itemize}
\item 第 1 轮：动作=summarize\_progress，阶段=materials，闭环状态=False。
\item 缺失字段：materials.volume\_material\_map、source.type、source.particle、source.energy、source.position、source.direction、physics.physics\_list、output.format、output.path、geometry.ask.stack.footprint、geometry.ask.stack.thicknesses、geometry.ask.stack.layer\_clearance、geometry.ask.stack.parent\_size、geometry.ask.stack.nest\_clearance。
\item LLM 使用=True；回退=无；后端=llm\_semantic\_frame+runtime\_semantic。
\item 第 2 轮：动作=confirm\_overwrite，阶段=source，闭环状态=False。
\item 缺失字段：无。
\item LLM 使用=True；回退=无；后端=llm\_slot\_frame+runtime\_semantic。
\item 第 3 轮：动作=finalize，阶段=physics，闭环状态=True。
\item 缺失字段：无。
\item LLM 使用=True；回退=无；后端=llm\_semantic\_frame+runtime\_semantic+confirmed\_pending\_overwrite。
\end{itemize}
\subsection{案例 MTM16（nested\_target）}
\textbf{对话：}
\begin{itemize}
\item 用户：我想在一个盒子里嵌一个圆柱靶，先别急着补全。
\item 系统：好的，我们已确认了您想要的几何结构类型（盒中嵌圆柱）。接下来还需要您提供以下信息，以便继续配置：

1.  **几何**：外层盒子的具体尺寸，以及内嵌圆柱的半径。
2.  **材料**：为盒体和圆柱体分别指定材料。
3.  **源**：设置粒子源的类型、粒子种类和能量。
4.  **物理**：选择模拟所需的物理过程列表。
5.  **输出**：设定输出数据的格式和保存路径。

请随时告诉我其中任何一项的具体内容。
\item 用户：外盒 120 mm x 120 mm x 120 mm，内嵌铜圆柱半径 20 mm、半长 40 mm，嵌套间隙 2 mm；材料 ?；束流 gamma 1 MeV，位置 (0,0,-150) mm，方向 +z；物理列表 QBBC；输出 json。
\item 系统：我注意到您新提供的几何参数会覆盖之前已确认的设置。具体来说，外层盒子的尺寸、内嵌圆柱的半径和半长，以及嵌套间隙都将被更新。

如果您确定要应用这些新数值，请直接回复“确认”；如果您希望保留原来的参数，请回复“保持原值”。
\item 用户：确认
\item 系统：好的，配置已完成。目前我们已确认的参数包括：嵌套间隙、内嵌圆柱半长、内嵌圆柱半径以及外层容器的尺寸。
\end{itemize}
\textbf{内部链路（自然语言）：}
\begin{itemize}
\item 第 1 轮：动作=summarize\_progress，阶段=materials，闭环状态=False。
\item 缺失字段：materials.volume\_material\_map、source.type、source.particle、source.energy、source.position、source.direction、physics.physics\_list、output.format、output.path、geometry.ask.nest.parent\_size、geometry.ask.nest.child\_radius、geometry.ask.nest.child\_half\_length、geometry.ask.nest.clearance。
\item LLM 使用=True；回退=无；后端=llm\_slot\_frame+runtime\_semantic。
\item 第 2 轮：动作=confirm\_overwrite，阶段=source，闭环状态=False。
\item 缺失字段：无。
\item LLM 使用=True；回退=无；后端=llm\_slot\_frame+runtime\_semantic。
\item 第 3 轮：动作=finalize，阶段=physics，闭环状态=True。
\item 缺失字段：无。
\item LLM 使用=True；回退=无；后端=llm\_semantic\_frame+runtime\_semantic+confirmed\_pending\_overwrite。
\end{itemize}
\subsection{案例 MTM17（coaxial\_shield）}
\textbf{对话：}
\begin{itemize}
\item 用户：I need a shell-style shielding toy model, but the shell thicknesses are still open.
\item 系统：Geometry: updated geometry structure; still needs shell inner radius, shell thicknesses. Materials: still needs volume-to-material binding. Source: updated source energy, source position, source direction; still needs particle type. Physics: updated physics list. Output: still needs output format, output path.
\item 用户：Use a shell with inner radius 15 mm, thicknesses 5 mm and 8 mm, half length 40 mm, nested core radius 10 mm, nested core half length 30 mm, clearance 1 mm, material lead; isotropic neutron source 2 MeV at (0,0,0) mm; physics Shielding; output hdf5.
\item 系统：This would overwrite an already confirmed value. The pending changes are:
- Source type: from "point" to "isotropic"
- Source energy: from 1.0 MeV to 2.0 MeV

Please reply with **'confirm'** to apply these changes, or **'keep existing'** to keep the current values.
\item 用户：confirm
\item 系统：Configuration complete. The following items are confirmed: source position, source energy, source type, and shell clearance.
\end{itemize}
\textbf{内部链路（自然语言）：}
\begin{itemize}
\item 第 1 轮：动作=summarize\_progress，阶段=materials，闭环状态=False。
\item 缺失字段：materials.volume\_material\_map、source.particle、output.format、output.path、geometry.ask.shell.inner\_radius、geometry.ask.shell.thicknesses、geometry.ask.shell.half\_length、geometry.ask.shell.child\_radius、geometry.ask.shell.child\_half\_length、geometry.ask.shell.clearance。
\item LLM 使用=True；回退=无；后端=llm\_slot\_frame+runtime\_semantic。
\item 第 2 轮：动作=confirm\_overwrite，阶段=source，闭环状态=False。
\item 缺失字段：无。
\item LLM 使用=True；回退=无；后端=llm\_slot\_frame+runtime\_semantic。
\item 第 3 轮：动作=finalize，阶段=physics，闭环状态=True。
\item 缺失字段：无。
\item LLM 使用=True；回退=无；后端=llm\_semantic\_frame+runtime\_semantic+confirmed\_pending\_overwrite。
\end{itemize}
\subsection{案例 MTM18（boolean\_cavity）}
\textbf{对话：}
\begin{itemize}
\item 用户：我想做一个挖空的盒体几何，但先只给你意图。
\item 系统：好的，目前进度如下：

**已设置**：几何结构类型（挖空盒体）。
**待补充**：为了完成挖空操作，还需要您提供两个部件的具体尺寸（部件A的x、y、z尺寸，以及部件B的x、y尺寸）。

此外，材料、粒子源、物理过程和输出设置等信息也仍待补充。我们一步步来。
\item 用户：做减法布尔体：A 为 40 mm x 40 mm x 40 mm 盒体，B 为 20 mm x 20 mm x 20 mm 盒体，从 A 中减去 B；材料 铝；点源 gamma 0.3 MeV，位置 (0,0,-60) mm，方向 +z；物理列表 QBBC；输出 xml。
\item 系统：我检测到这会覆盖已经确认的值。待修改内容：geometry.graph\_program：\{'nodes': [\{'id': 'a', 'type': 'Box', 'x': 29.287457362487178, 'y': 16.313688044337646, 'z': 53.82008547798903\}, \{'id': 'b', 'type': 'Box', 'x': 10.432721500065707, 'y': 45.19115032300169, 'z': 32.42666876844392\}, \{'id': 'bool', 'type': 'BooleanBinary', 'op': 'subtraction', 'left': 'a', 'right': 'b'\}], 'root': 'boolean', 'constraints': []\} -> \{'nodes': [\{'id': 'a', 'type': 'Box', 'x': 40.0, 'y': 40.0, 'z': 40.0\}, \{'id': 'b', 'type': 'Box', 'x': 20.0, 'y': 20.0, 'z': 20.0\}, \{'id': 'bool', 'type': 'BooleanBinary', 'op': 'subtraction', 'left': 'a', 'right': 'b'\}], 'root': 'boolean', 'constraints': []\}。如果要应用这些修改，请直接回复“确认”；如果要保留原值，请直接说“保持原值”。
\item 用户：确认
\item 系统：配置已完成。当前已确认的设置有：布尔 solid A 和 B 的尺寸、输出路径以及输出格式。
\end{itemize}
\textbf{内部链路（自然语言）：}
\begin{itemize}
\item 第 1 轮：动作=summarize\_progress，阶段=geometry，闭环状态=False。
\item 缺失字段：geometry.params.bool\_a\_x、geometry.params.bool\_a\_y、geometry.params.bool\_a\_z、geometry.params.bool\_b\_x、geometry.params.bool\_b\_y、geometry.params.bool\_b\_z、materials.volume\_material\_map、source.type、source.particle、source.energy、source.position、source.direction、physics.physics\_list、output.format、output.path、geometry.ask.boolean.solid\_a\_size、geometry.ask.boolean.solid\_b\_size。
\item LLM 使用=True；回退=无；后端=llm\_slot\_frame+runtime\_semantic。
\item 第 2 轮：动作=confirm\_overwrite，阶段=materials，闭环状态=False。
\item 缺失字段：无。
\item LLM 使用=True；回退=无；后端=llm\_slot\_frame+runtime\_semantic。
\item 第 3 轮：动作=finalize，阶段=source，闭环状态=True。
\item 缺失字段：无。
\item LLM 使用=True；回退=无；后端=llm\_semantic\_frame+runtime\_semantic+confirmed\_pending\_overwrite。
\end{itemize}
\subsection{案例 MTM19（ring\_complete）}
\textbf{对话：}
\begin{itemize}
\item 用户：Create a complete ring setup: 16 modules, each 6 mm x 6 mm x 2 mm, radius 32 mm, clearance 0.5 mm, material silicon; point gamma source 0.511 MeV at (0,0,0) mm; physics QBBC; output json.
\item 系统：Your configuration is complete. I have set up the ring geometry with 16 silicon modules, a 0.511 MeV gamma point source at the origin, and the QBBC physics list. Output will be in JSON format.

The following items are confirmed: volume-to-material binding, output path, clearance, and radius.
\end{itemize}
\textbf{内部链路（自然语言）：}
\begin{itemize}
\item 第 1 轮：动作=finalize，阶段=materials，闭环状态=True。
\item 缺失字段：无。
\item LLM 使用=True；回退=无；后端=llm\_slot\_frame+runtime\_semantic。
\end{itemize}
\subsection{案例 MTM20（grid\_complete）}
\textbf{对话：}
\begin{itemize}
\item 用户：请创建完整的阵列配置：3 x 3 盒体阵列，每个模块 12 mm x 12 mm x 3 mm，pitch\_x 15 mm，pitch\_y 15 mm，间隙 1 mm，材料 碘化铯；点源 gamma 0.662 MeV，位置 (0,0,-90) mm，方向 +z；物理列表 QBBC；输出 root。
\item 系统：配置已完成。当前已确认的设置包括：体积与材料绑定、输出路径、阵列间隙（clearance）、以及 y 方向和 x 方向的阵列间距（pitch）。
\end{itemize}
\textbf{内部链路（自然语言）：}
\begin{itemize}
\item 第 1 轮：动作=finalize，阶段=materials，闭环状态=True。
\item 缺失字段：无。
\item LLM 使用=True；回退=无；后端=llm\_slot\_frame+runtime\_semantic。
\end{itemize}
\subsection{案例 MTM21（layered\_detector）}
\textbf{对话：}
\begin{itemize}
\item 用户：I want a stacked detector slab, but the individual layer thicknesses are not fixed yet.
\item 系统：Here's what we have so far:

**Geometry:** You've set up a stacked detector structure. We still need the footprint (size) of each layer and their individual thicknesses.

**Materials:** We need to assign materials to the detector volumes.

**Source:** We need the source type, particle type, and its energy and position.

**Physics \& Output:** We still need to select the physics list and specify the output format and file path.
\item 用户：Use stack footprint 25 mm x 25 mm, thicknesses 1 mm, 2 mm, 1 mm, layer clearance 0.2 mm, outer box 30 mm x 30 mm x 12 mm, container clearance 0.5 mm, material silicon; beam electron 5 MeV from (0,0,-50) mm along +z; physics FTFP\_BERT; output json.
\item 系统：Your new geometry settings would overwrite previously confirmed values.

**Pending Change:**
- **Stack footprint:** 25 mm x 25 mm (was \textasciitilde{}16.3 mm x \textasciitilde{}10.3 mm)
- **Stack layer thicknesses:** 1 mm, 2 mm, 1 mm (was \textasciitilde{}6.68 mm, \textasciitilde{}1.19 mm, \textasciitilde{}5.59 mm)
- **Layer clearance:** 0.2 mm (was \textasciitilde{}0.73 mm)
- **Outer container size:** 30 mm x 30 mm x 12 mm (was \textasciitilde{}18.1 mm x \textasciitilde{}45.5 mm x \textasciitilde{}16.7 mm)
- **Container clearance:** 0.5 mm (was \textasciitilde{}0.87 mm)

Reply **'confirm'** to apply these new values, or **'keep existing'** to keep the current geometry.
\item 用户：confirm
\item 系统：Your configuration is complete. The following geometry items are confirmed: stack container clearance, stack container size, stack layer clearance, and stack layer thicknesses.
\end{itemize}
\textbf{内部链路（自然语言）：}
\begin{itemize}
\item 第 1 轮：动作=summarize\_progress，阶段=materials，闭环状态=False。
\item 缺失字段：materials.volume\_material\_map、source.type、source.particle、source.energy、source.position、source.direction、physics.physics\_list、output.format、output.path、geometry.ask.stack.footprint、geometry.ask.stack.thicknesses、geometry.ask.stack.layer\_clearance、geometry.ask.stack.parent\_size、geometry.ask.stack.nest\_clearance。
\item LLM 使用=True；回退=无；后端=llm\_semantic\_frame+runtime\_semantic。
\item 第 2 轮：动作=confirm\_overwrite，阶段=source，闭环状态=False。
\item 缺失字段：无。
\item LLM 使用=True；回退=无；后端=llm\_slot\_frame+runtime\_semantic。
\item 第 3 轮：动作=finalize，阶段=physics，闭环状态=True。
\item 缺失字段：无。
\item LLM 使用=True；回退=无；后端=llm\_semantic\_frame+runtime\_semantic+confirmed\_pending\_overwrite。
\end{itemize}
\subsection{案例 MTM22（embedded\_sensor）}
\textbf{对话：}
\begin{itemize}
\item 用户：请创建完整的嵌套配置：外盒 80 mm x 80 mm x 80 mm，内嵌铅圆柱半径 15 mm、半长 25 mm，嵌套间隙 1 mm；材料 ?；点源 gamma 1 MeV，位置 (0,0,-100) mm，方向 +z；物理列表 Shielding；输出 csv。
\item 系统：配置已完成。当前已确认的设置包括：体积与材料绑定、输出路径、嵌套间隙（clearance）以及圆柱半径。
\end{itemize}
\textbf{内部链路（自然语言）：}
\begin{itemize}
\item 第 1 轮：动作=finalize，阶段=materials，闭环状态=True。
\item 缺失字段：无。
\item LLM 使用=True；回退=无；后端=llm\_slot\_frame+runtime\_semantic。
\end{itemize}
\subsection{案例 MTM23（layered\_shield\_complete）}
\textbf{对话：}
\begin{itemize}
\item 用户：Create a complete shell configuration: inner radius 12 mm, thicknesses 4 mm and 6 mm, half length 35 mm, nested core radius 8 mm, nested core half length 25 mm, clearance 1 mm, material iron; isotropic gamma source 0.8 MeV at (0,0,0) mm; physics QBBC; output hdf5.
\item 系统：Your configuration is complete. The following items are confirmed: box dimensions (X, Y, Z) and volume-to-material bindings.
\end{itemize}
\textbf{内部链路（自然语言）：}
\begin{itemize}
\item 第 1 轮：动作=finalize，阶段=materials，闭环状态=True。
\item 缺失字段：无。
\item LLM 使用=True；回退=无；后端=llm\_slot\_frame+runtime\_semantic。
\end{itemize}
\subsection{案例 MTM24（boolean\_complete）}
\textbf{对话：}
\begin{itemize}
\item 用户：请创建完整布尔配置：做并集布尔体，A 为 30 mm x 30 mm x 20 mm 盒体，B 为 10 mm x 10 mm x 40 mm 盒体，与 A 做 union；材料 铝；束流 gamma 0.5 MeV，位置 (0,0,-70) mm，方向 +z；物理列表 QBBC；输出 json。
\item 系统：配置已完成。当前已确认的设置有：盒体在 x、y、z 方向的尺寸，以及体积与材料的绑定。
\end{itemize}
\textbf{内部链路（自然语言）：}
\begin{itemize}
\item 第 1 轮：动作=finalize，阶段=materials，闭环状态=True。
\item 缺失字段：无。
\item LLM 使用=True；回退=无；后端=llm\_slot\_frame+runtime\_semantic。
\end{itemize}
\end{document}